\section{Root-Cause Reward and Searched Verifier}
\label{sec:reward}

% [F3 content: reward + searched verifier. Source: rex/scoring.py, rex/harness_synth.py]
\paragraph{Graded reward.}
The scorer (\texttt{rex/scoring.py}) decomposes credit so that ``service is back'' is
necessary but not sufficient:
\begin{equation}
r = 0.30\,d + 0.25\,f + 0.45\,s - 0.60\,t,
\label{eq:reward}
\end{equation}
where \(d\) is diagnosis correctness (an LLM judge over the failure \emph{mechanism}, made
phrasing-robust), \(f\) is correct-fix credit, \(s\) is genuine resolution, and \(t\) is the
trap indicator. The trap penalty is what makes the loudest-alert decoy costly rather than
free, and the mechanism-level judge is the anti-reward-hacking lever. Empirically the reward
separates ``looks resolved'' (\(0.45\)) from a clean win (\(1.0\)) with the trap at
\(-0.60\).

\paragraph{Searched safety verifier.}
Rather than hand-writing the safety harness, we \emph{search} it. A Thompson-sampling tree
(\texttt{rex/tree.py}, \texttt{rex/harness\_synth.py}) explores rules-as-data --- declarative
predicates over the action and state, with \emph{no} LLM code execution --- and selects a
compact rule set that gates unsafe actions. The objective rewards held-out gating accuracy and
penalises rule count, yielding a small, auditable harness instead of an opaque classifier.

\paragraph{Why rules-as-data.}
Keeping the verifier declarative makes it inspectable, deterministic at eval time, and safe to
ship: a discovered rule can be read by an SRE before it is trusted to block a production
action.
